\section{study of integral functions}

\begin{theorem}[derivative of an integral function]
Let $I\subset \RR$ be an interval and $f\colon I \to \RR$ a locally
Riemann-integrable function. Let $a,b\colon A \to I$ be functions defined on a
set $A \subset \RR$. Then the function $F\colon A \to \RR$ is well-defined as
\[
  F(x) = \int_{a(x)}^{b(x)} f(t)\, dt.
\]
If furthermore $f$ is continuous and $a,b$ are differentiable, then $F$ is also
differentiable and we have
\[
  F'(x) = f(b(x)) \cdot b'(x) - f(a(x)) \cdot a'(x).
\]
\end{theorem}
%
\begin{proof}
By the stated hypotheses, for every $x\in A$, $f$ is integrable over the
interval $[a(x),b(x)]$, and thus the function $F$ is well-defined. Fixing
$x_0\in I$, we can consider the integral function
\[
  G(x) = \int_{x_0}^x f(t)\, dt
\]
and write
\[
  F(x)
    = \Enclose{G(x)}_{a(x)}^{b(x)}
    = G(b(x)) - G(a(x)).
\]
If $f$ is continuous, we can apply the fundamental theorem of calculus, which
ensures that $G$ is differentiable and $G'(x) = f(x)$ for every $x\in I$.
Therefore, if $a$ and $b$ are also differentiable, we have, by the chain rule:
\[
F'(x)
= G'(b(x)) \cdot b'(x) - G'(a(x)) \cdot a'(x)
= f(b(x)) \cdot b'(x) - f(a(x)) \cdot a'(x).
\]
\end{proof}

\begin{example}
The function
\[
 F(x) = \int_{x^2}^{x^4} \frac{1}{\ln t}\, dt
\]
is defined for every $x\in \RR\setminus\ENCLOSE{0,1,-1}$ since in this case the
interval $[x^2,x^4]$ is contained within the domain of the integrand function.
Furthermore, we have
\[
F'(x)
= \frac{4x^3}{\ln (x^4)} - \frac{2x}{\ln (x^2)}
= \frac{x^3-x}{\ln \abs{x}}.
\]
\end{example}

\begin{theorem}[integral of the small $o$]
Let $f(x)$ be a continuous function defined on an interval $[0,b]$ such that as
$x\to 0^+$, we have $f(x) = o(x^\alpha)$ for some $\alpha\ge 0$. Then, setting
\[
  F(x) = \int_0^x f(t)\, dt
\]
it follows that $F(x) = o(x^{\alpha+1})$. More concisely, we can write
\[
  \int_0^x o(t^\alpha)\, dt = o (x^{\alpha+1}),
  \qquad \text{as $x\to 0^+$}
\]
if the integrand function is continuous.
\end{theorem}
%
\begin{proof}
By the mean value theorem for integrals (theorem~\ref{th:media_integrale}), for
every $x\in [0,b]$, there must exist $c(x)\in[0,x]$ such that
\[
  F(x)
  = \int_0^x f(t)\, dt
  = x \cdot f(c(x)).
\]
Thus, since $0\le c(x)\le x$, we have
\[
\abs{\frac{F(x)}{x^{\alpha+1}}}
= \abs{\frac{f(c(x))}{x^\alpha}}
= \abs{\frac{f(c(x))}{c^\alpha(x)}}
\cdot \abs{\frac{c(x)}{x}}^\alpha
\le \abs{\frac{f(c(x))}{c^\alpha(x)}} \to 0
\]
since $f(x) = o(x^\alpha)$ and as $x\to 0$, also $c(x)\to 0$, we indeed have
\[
  \lim_{x\to 0} \frac{f(c(x))}{c^\alpha(x)} = 0.
\]
\end{proof}

\begin{example}
We wish to calculate
\[
  \lim_{x\to 0^+} \frac{1}{x^4}\int_{\sin^2 x}^{\sin x} \frac{2- t\sin t - 2 \cos t}{e^t - 1}\, dt.
\]
Calling $f(x)$ the integrand function, it is not difficult to verify, using
Taylor's formulas, that
\[
  f(x)
  = \frac{\frac{x^4}{12}-o(x^4)}{x+o(x)}
  = \frac{x^3}{12} + o(x^3), \qquad \text{as $x\to 0$}.
\]
In particular, the function $f$ can be extended by continuity by setting
$f(0)=0$. Thus, by the previous theorem,
\[
  F(x) = \int_0^x f(t)\, dt
  = \int_0^x \enclose{\frac{t^3}{12}+o(t^3)}\, dt
  = \frac{x^4}{48} + o(x^4)
\]
from which
\begin{align*}
 \frac{1}{x^4} \int_{\sin^2 x}^{\sin x}
 f(t) \, dt
 &= \frac{F(\sin x) - F(\sin^2 x)}{x^4}\\
  &= \frac{\frac{(\sin x)^4}{48} + o(\sin^4 x) - \frac{(\sin^2 x)^4}{48} +
 o(\sin^8 x)}{x^4} \\
 &= \frac{\frac{x^4}{48} + o(x^4)}{x^4} \to \frac{1}{48}.
\end{align*}
\end{example}

\begin{example}
We wish to calculate
\[
  \lim_{x\to 0^+} \int_{x}^{2x} \frac{1}{t+\sin t}\, dt.
\]
Setting
\[
  f(x) = \frac{1}{x+\sin x}
\]
we have
\[
 f(x)
  = \frac{1}{2x+o(x^2)} = \frac{1}{2x(1+o(x))}
  = \frac{1+o(x)}{2x} = \frac{1}{2x} + o(1).
\]
The function $f(x)-\frac 1{2x} = o(1)$ can be extended by continuity even at
$x=0$, and thus we can apply the previous theorem to obtain
\begin{align*}
  \int_x^{2x} f(t)\, dt
  &= \int_x^{2x} \enclose{\frac 1 {2t} + o(1)}\, dt \\
  &= \frac 1 2 \Enclose{\ln t}_x^{2x} +  \Enclose{o(t)}_x^{2x} \\
  &= \frac {\ln 2x - \ln x} 2 + o(2x) - o(x)
  = \frac {\ln 2} 2 + o(x) \to \frac{\ln 2}{2}.
\end{align*}
\end{example}

\begin{theorem}
Let $a,b\colon J \to \RR$ be two functions defined on an interval $J$ with
values in an interval $I$. Let $x_0$ be an accumulation point of $J$ and let $a$
be an accumulation point of $I$. Suppose further that as $x\to x_0$, we have
$a(x)\to a$, $b(x)\to a$. Let $f,g\colon I \to \RR$ be continuous functions. If
$f(x)\sim g(x)$ as $x\to a$, then we have
\[
  \int_{a(x)}^{b(x)} f(t)\, dt
  \sim \int_{a(x)}^{b(x)} g(t)\, dt
  \qquad \text{as $x\to x_0$}.
\]
\end{theorem}
%
\begin{proof}
Let $F$ be an antiderivative of $f$ and $G$ an antiderivative of $g$. Then
\[
  \frac{\displaystyle\int_{a(x)}^{b(x)}f(t)\, dt}{\displaystyle\int_{a(x)}^{b(x)} g(t)\, dt}
  = \frac{F(b(x)) - F(a(x))}{G(b(x))-G(a(x))}.
\]
Applying Cauchy's theorem~\ref{th:cauchy}, we have
\[
\frac{F(b(x)) - F(a(x))}{G(b(x))-G(a(x))}
= \frac{F'(c(x))}{G'(c(x))}
= \frac{f(c(x))}{g(c(x))}
\]
for some $c(x)$ between $a(x)$ and $b(x)$. Since $a(x)\to a$ and $b(x)\to a$, it
follows that $c(x)\to a$ as $x\to x_0$. Furthermore, since $f(x)\sim g(x)$ as
$x\to a$, by changing the variable in the limit, we can deduce that
\[
\frac{f(c(x))}{g(c(x))} \to 1 \qquad \text{as $x\to x_0$}.
\]
\end{proof}

\begin{example}
We wish to calculate
\[
  \lim_{x\to 0 }\int_{x^2}^{3x^2} \frac{\tg t}{t^2}\, dt.
\]
Since as $x\to 0$ we have
\[
  \frac{\tg x}{x^2} \sim \frac 1 x
\]
and since
\[
  \int_{x^2}^{3x^2}\frac{1}{t}\, dt
  = \ln(3x^2) - \ln(x^2) = \ln 3
\]
thanks to the previous theorem, we can deduce that the desired limit is indeed
$\ln 3$.
\end{example}